% Source: https://artofproblemsolving.com/community/c3h1266946_halllp_what_did_i_do_wrong_wallbash
% Stable identity: aops-topic-1266946-q01
All interior angles are $120^\circ$. Extending the first, third and fifth sides (lengths $2$, $6$, $4$) creates three equilateral corner triangles, on the sides of lengths $4$, $2$ and $6$, and the hexagon together with them forms an equilateral triangle.

\begin{asy}
usepackage("amsmath");
size(5.5cm);
defaultpen(fontsize(10pt));
pair[] P; P[0]=(0,0); real[] s={2,4,6,2,4,6};
for(int i=0;i<6;++i) P[i+1]=P[i]+s[i]*dir(60*i);
pair T1=P[0]+(s[0]+s[1]+s[2])*dir(0)-(0,0);
draw(P[0]--P[1]--P[2]--P[3]--P[4]--P[5]--cycle,linewidth(1));
pair X=extension(P[0],P[1],P[2],P[3]), Y=extension(P[2],P[3],P[4],P[5]), Z=extension(P[4],P[5],P[0],P[1]);
draw(P[1]--X--P[2],dashed); draw(P[3]--Y--P[4],dashed); draw(P[5]--Z--P[0],dashed);
for(int i=0;i<6;++i) label("$"+string(s[i])+"$",(P[i]+P[i+1])/2,unit((P[i]+P[i+1])/2-(P[0]+P[1]+P[2]+P[3]+P[4]+P[5])/6));
label("\mbox{}",(min(currentpicture,user=true)+max(currentpicture,user=true))/2,black); // keeps the text after the figure black when LaTeX typesets the labels inline
\end{asy}

The big triangle's side is $4+6+2=12$ (one hexagon side plus the two corner-triangle sides on it). The corner triangles have sides $4$, $2$ and $6$. Hence
\[\text{Area}=\dfrac{\sqrt3}{4}\left(12^2-4^2-2^2-6^2\right)=\dfrac{\sqrt3}{4}\cdot 88=\boxed{22\sqrt3}.\]
